\begin{abstract}
A vision-language model that has genuinely unified its modalities should treat a
piece of content as the same content whether it arrives as text or as pixels.
We make that intuition measurable. Our central obstacle is that cross-modal
equivalence is ill-posed as usually stated: a phrase and an image region are not
the same object, occupy different numbers of tokens, and carry different
information, so any raw distance between their representations confounds
representational divergence with genuine content mismatch. We resolve this with
three constructions. First, equivalence is asserted not on the substituted span
but at the \emph{merge position} --- the first token of a shared, verbatim-identical
suffix --- where the two views must have reconciled whatever they encoded.
Second, we normalise the cross-modal distance by a \emph{within-modality} control
distance measured on the same items, yielding the Modality Substitution Gap
($\msg$), a scale-free ratio whose value of $1$ means crossing modalities costs
no more than re-expressing content within one. Third, we introduce
\emph{orthographic substitution} --- replacing a word with a picture of that word ---
which makes the substitution exactly content-preserving and so removes the
information asymmetry that makes referential grounding an unreliable probe.
Around these we build the controls the measurement needs but that prior modality-gap
work does not report: a span-free floor establishing that a comprehension probe
depends on the substituted content at all, and a hierarchy of nulls
(translation, rotation, arbitrary linear map, each fitted out-of-fold) that
distinguishes a gap in \emph{information} from a gap in \emph{coordinates}. Applied to a
2B-parameter VLM, the instrument produces a large and reproducible gap under
orthographic substitution, and the accompanying controls determine whether that
gap reflects anything a downstream computation could not undo.
\emph{(Numerical results are pending the final run; Section~5 is a placeholder.)}
\end{abstract}
